\section{Addressing}

Because the layer growth follows a fixed recurrence, every dock has a short canonical name, a string of digits, and a step to a neighbor is arithmetic on that string. Finding any dock among astronomically many costs about the logarithm of their number.

This matters twice over. It makes the mesh a searchable space, not just a place, so the simulator can build and navigate large patches of the exact mesh, and so the mesh can serve as a computer with a solved address for every cell. And it gives the bulk a natural memory: a stored pattern hangs at an address on the addressing tree, still if it is a remembered thing and active if it is a running program, recalled by navigating to the address.

The same exponential branching that makes the bulk roomy makes the addressing tree deep and wide, so the tree of names is at once a trie, a content-addressable memory, a balanced logarithmic index, and an error-correcting code. All of this needs the exponential room that only hyperbolic space provides. So the geometry that gives holography and the mass hierarchy also gives the mesh a fast, capacious, self-organizing place to store and find structure, which is the bridge from physics to computation taken up later.
